% Part: normal-modal-logic
% Chapter: frame-definability
% Section: first-order-definability

\documentclass[../../../include/open-logic-section]{subfiles}

\begin{document}

\olfileid{nml}{frd}{fol}

\olsection{પ્રથમ-ક્રમ વ્યાખ્યાયિતતા}

આપણે જોયું છે કે ફ્રેમોના સુલભતા સંબંધના કેટલાંક
ગુણધર્મો મોડલ !!{formula}sથી વ્યાખ્યાયિત કરી શકાય છે.
દાખલા તરીકે, ફ્રેમોની સંમિતતા !!{formula}~\Ax{B},
$p \lif \Box\Diamond
p$થી વ્યાખ્યાયિત થાય છે. અત્યાર સુધી મળેલી બધી
શરતો એક દ્વિસ્થાની !!{predicate} ધરાવતી
પ્રથમ-ક્રમની !!a{language}માંનાં !!{formula}sથી
વ્યક્ત કરી શકાય છે. દાખલા તરીકે, સંમિતતા
$\lforall[x][\lforall[y][(\Atom{Q}{x,y} \lif \Atom{Q}{y, x})]]$થી
એ અર્થમાં વ્યાખ્યાયિત થાય છે કે $\Domain{M} =
W$ અને $\Assign{Q}{M} = R$ ધરાવતી પ્રથમ-ક્રમની
!!{structure}~$\Struct{M}$ એ ઉપરનું સૂત્ર ત્યારે અને
માત્ર ત્યારે જ સંતોષે છે જ્યારે $R$ સંમિત હોય.
આથી નીચેની વ્યાખ્યા સૂચાય છે:

\begin{defn}
  ફ્રેમોનો વર્ગ~$\mClass{F}$ \emph{પ્રથમ-ક્રમમાં વ્યાખ્યાયિત}
  છે, જો એક દ્વિસ્થાની !!{predicate}~$Q$ ધરાવતી
  પ્રથમ-ક્રમની ભાષામાં એવું !!a{sentence}~$!A$
  હોય કે $\mModel{F} =
  \tuple{W, R} \in \mClass{F}$ ત્યારે અને માત્ર ત્યારે જ
  જ્યારે $\Domain{M} = W$ અને $\Assign{Q}{M}
  = R$ ધરાવતી પ્રથમ-ક્રમની !!{structure}~$\Struct{M}$માં
  $\Sat{M}{!A}$ હોય.
\end{defn}

અત્યાર સુધી વિચારેલા ગુણધર્મો અને તેમને વ્યાખ્યાયિત
કરતાં મોડલ !!{formula}s અપવાદરૂપ છે. દરેક
!!{formula} ફ્રેમોના પ્રથમ-ક્રમમાં વ્યાખ્યાયિત વર્ગને
વ્યાખ્યાયિત કરતું નથી; અને ફ્રેમોના દરેક પ્રથમ-ક્રમમાં
વ્યાખ્યાયિત વર્ગને કોઈ મોડલ સૂત્રથી વ્યાખ્યાયિત
કરી શકાતો નથી.

પહેલા દાવાનું વિરોધી દૃષ્ટાંત લોબનું સૂત્ર છે:
\begin{equation}
\Box(\Box p \lif p) \lif \Box p. \tag{\Ax{W}}
\end{equation}
\Ax{W} એ પરંપરિત અને પ્રતિલોમ-સુસ્થાપિત ફ્રેમોનો
વર્ગ વ્યાખ્યાયિત કરે છે. સંબંધ સુસ્થાપિત હોય છે જો
$Rw_2w_1$, $Rw_3w_2$, \dots{} ધરાવતી અનંત શ્રેણી
$w_1$, $w_2$, \dots{} ન હોય. દાખલા તરીકે,
~$\Nat$ પરનો $<$ સંબંધ સુસ્થાપિત છે, જ્યારે
~$\Int$ પરનો $<$ સંબંધ એવો નથી. સંબંધ
પ્રતિલોમ-સુસ્થાપિત ત્યારે અને માત્ર ત્યારે જ છે
જ્યારે તેનો પ્રતિલોમ સુસ્થાપિત હોય. તેથી
પ્રતિલોમ-સુસ્થાપિત સંબંધોમાં $Rw_1w_2$,
$Rw_2w_3$, \dots{} ધરાવતી અનંત શ્રેણી
$w_1$, $w_2$, \dots{} હોતી નથી.

પરંતુ પરંપરિત પ્રતિલોમ-સુસ્થાપિત સંબંધોને
વ્યાખ્યાયિત કરતું કોઈ પ્રથમ-ક્રમનું !!{formula}
નથી. કેમ કે માનો કે~$\Sat{M}{!F}$ ત્યારે અને માત્ર
ત્યારે જ જ્યારે $R =
\Assign{Q}{M}$ પરંપરિત અને પ્રતિલોમ-સુસ્થાપિત હોય.
સઘનતાની દલીલ માટે ભાષામાં નવા અચલ સંકેતો
$a_1$, $a_2$, \dots{} ઉમેરો. હવે $!A_n$ નીચેનું
!!{formula} હોય, જ્યાં સૂચક ઓછામાં ઓછો બે છે:
\[
(\Atom{Q}{a_1,a_2} \land \dots \land
    \Atom{Q}{a_{n-1},a_{n}})
\]
હવે !!{formula}sનો ગણ લો:
\[
\Gamma = \{!F, !A_2, !A_3, \dots\}.
\]
$\Gamma$નો દરેક સાન્ત ઉપગણ સંતોષ્ય છે: ઉપગણમાં
કોઈ $!A_k$ હોય તો તેમાંનું સૌથી મોટું સૂચક $k$
લો; એવું એકેય ન હોય તો બે લો. પછી
$\Domain{M_k} = \{1, \dots, k\}$,
$\Assign{a_i}{M_k} = i$ અને $\Assign{Q}{M_k} = <$ મૂકો.
$\{1,
\dots, k\}$ પરનો $<$ સંબંધ પરંપરિત અને
પ્રતિલોમ-સુસ્થાપિત હોવાથી $\Sat{M_k}{!F}$.
રચના મુજબ બે કે વધુ હોય એવા બધા $i \le k$ માટે
$\Sat{M_k}{!A_i}$ પણ છે.
પ્રથમ-ક્રમ તર્કશાસ્ત્રના સઘનતા પ્રમેયથી $\Gamma$
કોઈ !!{structure}~$\Struct{M}$માં સંતોષ્ય છે.
ધારણા મુજબ $\Sat{M}{!F}$ હોવાથી
$\Assign{Q}{M}$ પ્રતિલોમ-સુસ્થાપિત છે. પરંતુ
$\Assign{a_1}{M}$, $\Assign{a_2}{M}$,
\dots{} સ્પષ્ટ રીતે એવી અનંત શ્રેણી બનાવે છે
જે પ્રતિલોમ-સુસ્થાપિતતા પ્રતિબંધિત કરે છે.
\sourcecorrection{OLMD-017}{સ્થિર અંગ્રેજી મૂળમાં
શ્રેણી-સૂત્રનું $!A_1$ સૂચક અનિર્ધારિત છે: તેનું
પ્રથમ પગલું બીજા અચલથી શરૂ થાય છે. ગણને
બીજા સૂચકથી શરૂ કર્યો અને કોઈ શ્રેણી-સૂત્ર
વિના સાન્ત ઉપગણનો કિસ્સો પણ સ્પષ્ટ કર્યો છે.}
\sourcecorrection{OLMD-018}{સ્થિર અંગ્રેજી મૂળની
વ્યાખ્યામાં માત્ર દ્વિસ્થાની વિધેયચિહ્નવાળી ભાષા છે,
પરંતુ સઘનતાના પ્રમાણમાં નવા અચલ સંકેતો વપરાય છે.
પ્રમાણમાં ભાષાનો આ અચલ-વિસ્તાર સ્પષ્ટ કર્યો છે.}

બીજા દાવાનું વિરોધી દૃષ્ટાંત સાર્વત્રિકતાનો ગુણધર્મ
છે: દરેક $u$ અને $v$ માટે $Ruv$. સાર્વત્રિક
ફ્રેમો પ્રથમ-ક્રમના સૂત્ર
$\lforall[x][\lforall[y][\Atom{Q}{x,y}]]$થી
વ્યાખ્યાયિત થાય છે. પરંતુ એવો કોઈ મોડલ સૂત્ર
નથી જે બધી અને માત્ર સાર્વત્રિક ફ્રેમોમાં માન્ય હોય.
આ સ્વતંત્ર રીતે રસપ્રદ પરિણામનું ફળ છે: સાર્વત્રિક
ફ્રેમોમાં માન્ય સૂત્રો બરાબર સ્વવાચક, સંમિત અને
પરંપરિત ફ્રેમોમાં માન્ય સૂત્રો જ છે. સ્વવાચક,
સંમિત અને પરંપરિત છતાં અસાર્વત્રિક ફ્રેમો છે;
તેથી બધી સાર્વત્રિક ફ્રેમોમાં માન્ય દરેક !!{formula}
કેટલીક અસાર્વત્રિક ફ્રેમોમાં પણ માન્ય હોય છે.

\end{document}
